\documentclass[11pt,a4paper]{article}
\usepackage[utf8]{inputenc}
\usepackage[T1]{fontenc}
\usepackage[spanish]{babel}
\usepackage[margin=2.5cm]{geometry}
\usepackage{booktabs}
\usepackage{array}
\usepackage{listings}
\usepackage{xcolor}
\usepackage{enumitem}
\usepackage{parskip}
\usepackage{fancyhdr}
\usepackage{amsmath}
\usepackage{amssymb}

\definecolor{codegray}{RGB}{245,245,245}

\lstdefinestyle{log}{
  basicstyle=\ttfamily\small,
  backgroundcolor=\color{codegray},
  frame=single,
  framerule=0.4pt,
  xleftmargin=8pt,
  numbers=none,
  breaklines=true
}

\pagestyle{fancy}
\fancyhf{}
\rhead{\textcolor{gray}{\small TD7: Ingeniería de Datos — UTDT}}
\lhead{\textcolor{gray}{\small Examen Modelo 2 — Bloques 1, 2 y 4 (Papel)}}
\rfoot{\thepage}
\renewcommand{\headrulewidth}{0.4pt}

\setlist[enumerate,1]{label=\textbf{\alph*)}}
\setlist[enumerate,2]{label=\roman*.}

\begin{document}

\begin{center}
  {\LARGE\bfseries TD7 — Ingeniería de Datos}\\[6pt]
  {\large Universidad Torcuato Di Tella}\\[4pt]
  {\Large\bfseries Examen — Modelo 2}\\[2pt]
  {\large\bfseries Bloques 1, 2 y 4 \textnormal{---} \textit{Parte en papel}}\\[4pt]
  {\normalsize Duración: 2 horas \quad|\quad Dejar las máquinas de lado}
\end{center}
\noindent\rule{\linewidth}{0.8pt}
\vspace{2pt}

\noindent\textbf{Instrucciones:} Responda todos los ejercicios en papel. Justifique cada respuesta. No está permitido el uso de material de consulta ni de dispositivos electrónicos. Indique claramente el número de bloque y sub-ítem.

\vspace{4pt}

\begin{center}\small
\begin{tabular}{lll}
\toprule
\textbf{Bloque} & \textbf{Puntaje} & \textbf{Aprueba con} \\
\midrule
Bloque 1 — DER                    & 10 puntos & $\geq 5 / 10$ \\
Bloque 2 — Normalización + Concurrencia & 4 puntos  & $\geq 2 / 4$ \;(equiv.\ a 5/10) \\
Bloque 4 — Preguntas Teóricas     & 3 puntos  & $\geq 1{,}5 / 3$ \;(equiv.\ a 5/10) \\
\bottomrule
\end{tabular}
\end{center}

\vspace{4pt}
\noindent\rule{\linewidth}{0.4pt}

% ═══════════════════════════════════════════════════════════════
\section*{Bloque 1 — Diseño de DER \hfill\normalfont(10 puntos)}

Una plataforma de e-commerce desea modelar su sistema. Las reglas de negocio son:

\begin{itemize}
  \item Los \textbf{vendedores} tienen un identificador único, nombre, email y fecha de registro en la plataforma.
  \item Los \textbf{productos} tienen un código, nombre, descripción y precio. Cada producto pertenece a exactamente una \textbf{categoría}. Una categoría tiene un código, nombre y descripción. Una categoría puede existir aunque no tenga productos asignados.
  \item Cada producto es publicado por \emph{exactamente un} vendedor; un vendedor puede publicar varios productos o ninguno.
  \item Los \textbf{compradores} se identifican por un ID y tienen nombre, email y dirección de envío.
  \item Un comprador puede realizar uno o más \textbf{pedidos}. Cada pedido tiene un número único, fecha y estado (\texttt{pendiente}, \texttt{despachado}, \texttt{entregado}). Un pedido incluye uno o más productos con una cantidad determinada y el precio vigente al momento de la compra (puede diferir del precio actual del producto).
  \item Un comprador puede escribir \textbf{reseñas} de productos que haya comprado. Cada reseña tiene un puntaje (entero de 1 a 5) y un texto opcional. Un comprador puede reseñar un mismo producto como máximo una vez.
\end{itemize}

\begin{enumerate}
  \item (7 pts) Dibuje el DER completo. Indique entidades, atributos (marcando el identificador de cada entidad), interrelaciones, cardinalidades y tipos de participación (total o parcial).

  \vspace{12cm}

  \item (3 pts) Traduzca el DER al modelo relacional. Para cada tabla indique nombre, todos sus atributos, \underline{clave primaria} (subrayada) y \textit{claves foráneas} (en cursiva, indicando la tabla referenciada). Preste especial atención a la relación de pedido-producto (con atributos propios).
\end{enumerate}

\newpage

% ═══════════════════════════════════════════════════════════════
\section*{Bloque 2 — Normalización y Concurrencia \hfill\normalfont(4 puntos)}

% ─────────────────────────────────────────────────────────────
\subsection*{Ejercicio 2A — Normalización \hfill\normalfont(1{,}5 puntos)}

Se tiene la siguiente relación sin normalizar:

\begin{center}
\textbf{VENTA}(\texttt{nro\_venta, cod\_producto, fecha, cod\_cliente, nombre\_cliente,}\\
\texttt{descripcion\_producto, precio\_unitario, cantidad})
\end{center}

Las dependencias funcionales (DF) válidas son:

\begin{itemize}
  \item $\texttt{nro\_venta} \rightarrow \texttt{fecha},\; \texttt{cod\_cliente}$
  \item $\texttt{cod\_cliente} \rightarrow \texttt{nombre\_cliente}$
  \item $\texttt{cod\_producto} \rightarrow \texttt{descripcion\_producto},\; \texttt{precio\_unitario}$
  \item $\{\texttt{nro\_venta},\; \texttt{cod\_producto}\} \rightarrow \texttt{cantidad}$
\end{itemize}

\begin{enumerate}
  \item (0{,}5 pts) Identificar la clave primaria de \textbf{VENTA}. ¿Está la relación en 2FN? Justifique identificando las dependencias parciales si las hubiera.

  \vspace{3.5cm}

  \item (1 pt) Lleve la relación a \textbf{3FN}. Muestre todas las relaciones resultantes con sus atributos, clave primaria y claves foráneas. Verifique también que no queden dependencias transitivas.

  \vspace{5.5cm}
\end{enumerate}

% ─────────────────────────────────────────────────────────────
\subsection*{Ejercicio 2B — Concurrencia y Recuperabilidad \hfill\normalfont(2{,}5 puntos)}

Se tienen tres transacciones con las siguientes operaciones:

\begin{center}
$T_1:\; R(A),\; R(B),\; W(A)$ \qquad
$T_2:\; R(B),\; W(B)$ \qquad
$T_3:\; R(A),\; W(C)$
\end{center}

El siguiente schedule $S$ intercala sus operaciones:

\[
S:\; R_1(A),\; R_2(B),\; R_3(A),\; W_2(B),\; R_1(B),\; W_1(A),\; C_2,\; W_3(C),\; C_1,\; C_3
\]

\begin{enumerate}
  \item (0{,}75 pts) Identificar \textbf{todos} los pares de operaciones en conflicto en $S$. Construir el \textbf{grafo de precedencias}. ¿Es $S$ serializable por conflictos? Si lo es, indique el orden serial equivalente.

  \vspace{5cm}

  \item (0{,}75 pts) ¿El schedule $S$ es \textbf{recuperable}? ¿Existe riesgo de \textbf{rollback en cascada}? Analice quién leyó datos escritos por quién y en qué orden se produjeron los commits.

  \vspace{4cm}

  \item (1 pt) El sistema usa la estrategia \textbf{REDO} (actualización diferida). Al momento de la falla, el log en disco contiene:

\begin{lstlisting}[style=log]
(BEGIN,  T1)
(BEGIN,  T2)
(WRITE,  T2, B, 200)
(COMMIT, T2)
(WRITE,  T1, A, 500)
         <-- FALLA DEL SISTEMA AQUI
\end{lstlisting}

  \noindent\small\textit{Formato de WRITE en REDO: (WRITE, $T_i$, ítem, $v_{new}$). El valor anterior no se almacena en el log.}

  \normalsize\vspace{4pt}
  ¿Qué transacciones deben rehacerse y cuáles no? Describa paso a paso el proceso de recuperación e indique los valores finales de $A$ y $B$ en disco.

  \vspace{4.5cm}
\end{enumerate}

\newpage

% ═══════════════════════════════════════════════════════════════
\section*{Bloque 4 — Preguntas Teóricas \hfill\normalfont(3 puntos — 1 punto cada una)}

\begin{enumerate}

  \item \textbf{Almacenamiento — RAID.} Describir el funcionamiento de \textbf{RAID 0}, \textbf{RAID 1} y \textbf{RAID 5}, indicando para cada nivel sus ventajas y desventajas en términos de rendimiento, capacidad útil y tolerancia a fallos. ¿Cuál elegiría para un servidor de base de datos en producción donde la disponibilidad y la integridad de los datos son críticas? Justifique su elección.

  \vspace{6cm}

  \item \textbf{Índices — Selectividad.} ¿En qué situaciones \textbf{no} conviene crear un índice sobre un atributo y puede resultar más lento que un full scan? Relacione su respuesta con el concepto de \textbf{selectividad}. ¿Qué tipos de consultas se benefician más del uso de índices?

  \vspace{6cm}

  \item \textbf{Optimización de consultas — Estadísticas.} ¿Qué son las \textbf{estadísticas} de una tabla (por ejemplo, \texttt{n\_distinct}, histogramas de distribución de valores)? ¿Cómo las usa el optimizador para estimar el costo de un plan de ejecución? ¿Qué puede ocurrir si las estadísticas están desactualizadas?

  \vspace{6cm}

\end{enumerate}

\end{document}
